\begin{center}
\textbf{Aprendizagem automática explicável e reprodutível para o prognóstico da ELA: classificação da progressão funcional e previsão da sobrevivência}
\end{center}

\section*{Resumo}
\label{sec:resumo}

A esclerose lateral amiotrófica apresenta evolução variável, dificultando a previsão do declínio funcional e da sobrevivência. Esta dissertação utiliza a PRO-ACT em dois estudos prognósticos. O Estudo~I previu a progressão funcional rápida aos três e seis meses a partir de dados clínicos basais. No teste, o XGBoost obteve PR-AUC de 0,456 e identificou 72 dos 84 progressores rápidos, com sensibilidade de 85,7\% e precisão de 35,1\%. A análise SHAP destacou o ALSFRS-R basal e a capacidade vital forçada. O Estudo~II analisou a morte nos 24 meses após o início dos sintomas. O LightGBM com sobreamostragem aleatória obteve ROC-AUC de 0,923 e identificou 30 dos 35 doentes com sobrevivência curta. Em ambos, a ordenação não determinou o comportamento no limiar escolhido. A definição do resultado, o desequilíbrio, a separação e o limiar influenciaram a interpretação. É necessária validação externa e prospetiva antes do uso clínico.

\section*{Palavras-chave}
\label{sec:palavras}

Esclerose lateral amiotrófica, progressão funcional, prognóstico de sobrevivência, desequilíbrio de classes, seleção do limiar.
